% ----------------------------------------------------------
% Glossário
% ----------------------------------------------------------
%
% Consulte o manual da classe abntex2 para orientações sobre o glossário.
%
%\glossary

% ----------------------------------------------------------
% Apêndices
% ----------------------------------------------------------

% ---
% Inicia os apêndices
% ---

\begin{apendicesenv}

% Imprime uma página indicando o início dos apêndices
\partapendices

% ======================================================================
\chapter{Exemplos de aplicações GMTP}
\label{ap:aplicacoes}

Neste apêndice, apresentam-se alguns aplicações exemplo que utilizam o GMTP como protocolo de transporte.

\section{Constantes para uso do GMTP}
\label{ap:gmtp.h}

No código fonte a seguir, está disponível um arquivo de cabeçalho escrito na linguagem C, denominado $gmtp.h$, onde definem-se as principais constantes necessárias para s instanciar um \textit{socket} GMTP.

\begin{lstlisting}[language=C++, caption={Constantes e definições para uso das aplicações GMTP}, label={list:gmtp.h}]

#ifndef GMTP_H_
#define GMTP_H_

#define SOCK_GMTP     7
#define IPPROTO_GMTP  254
#define SOL_GMTP      281

enum gmtp_sockopt_codes {
	GMTP_SOCKOPT_FLOWNAME = 1,
	GMTP_SOCKOPT_MEDIA_RATE,
	GMTP_SOCKOPT_MAX_TX_RATE,
	GMTP_SOCKOPT_UCC_TX_RATE,
	GMTP_SOCKOPT_GET_CUR_MSS,
	GMTP_SOCKOPT_SERVER_RTT,
	GMTP_SOCKOPT_SERVER_TIMEWAIT,
	GMTP_SOCKOPT_PULL,
	GMTP_SOCKOPT_ROLE_RELAY,
	GMTP_SOCKOPT_RELAY_ENABLED,
	GMTP_SOCKOPT_UCC_TYPE
};

#endif /* GMTP_H_ */
\end{lstlisting}

No código definem-se o valor correspondente ao tipo do \textit{socket} do GMTP como $7$, o número do protocolo GMTP como $254$ (a IANA definiu o valor 254 para testes e experimentos) e o código do GMTP para ler ou alterar as opções do \textit{socket} GMTP como $281$.
Além disso, a enumeração \textit{gmtp\_sockopt\_codes} serve como identificador das opções de \textit{socket} do GMTP, na chamada das funções \textit{setsockopt} e \textit{getsockopt}.

\section{Cliente GMTP}
\label{ap:cliente-gmtp}

No código fonte a seguir, está disponível o código de implementação de um cliente GMTP escrito na linguagem C. No código, faz-se uso do arquivo $gmtp.h$, definido no Apêndice \ref{ap:gmtp.h}.

\begin{lstlisting}[language=C++, caption={Trecho de código de um cliente GMTP escrito na linguagem C.}, label={list:client-gmtp}]
#include <sys/socket.h>
#include <arpa/inet.h>
#include "gmtp.h"

#define SERVER_PORT 2000

int main(int argc, char **argv) {
	char *server_addr = argv[1];
	struct sockaddr_in addr;
	int sockfd = socket(AF_INET, SOCK_GMTP, IPPROTO_GMTP);
	
	addr.sin_family = AF_INET;
	addr.sin_addr.s_addr = inet_addr(server_addr);
	addr.sin_port = htons (SERVER_PORT);
	int ret = connect(sockfd, (struct sockaddr *)&addr,
						sizeof(addr));
						
	/* A partir deste ponto, a aplicacao pode receber dados
	 * atraves de funcoes como recv() ou recvmsg().
	 */
	
	return 0;
}
\end{lstlisting}

No exemplo, a porta do servidor é definida como $2000$ e o endereço IP do servidor deve ser passado como parâmetro de linha de comando da aplicação.

\section{Servidor GMTP}
\label{ap:servidor-gmtp}

No código fonte a seguir, está disponível o código de implementação de um servidor GMTP escrito na linguagem C. No código, faz-se uso do arquivo $gmtp.h$, definido no Apêndice \ref{ap:gmtp.h}.

\begin{lstlisting}[language=C++, caption={Trecho de código de um servidor GMTP escrito na linguagem C.}, label={list:servidor-gmtp}]
#include <sys/socket.h>
#include <arpa/inet.h>
#include "gmtp.h"

#define SERVER_PORT 2000

int main(int argc, char *argv[])
{
	int sockfd, newSocket;
	struct sockaddr_in addr;
	struct sockaddr_storage addr_s;
	int bitrate = 250000; /* 250000 B/s, ou 2Mbps */

	sockfd = socket(PF_INET, SOCK_GMTP, IPPROTO_GMTP);
	setsockopt(sockfd, SOL_GMTP, GMTP_SOCKOPT_MEDIA_RATE, &media_rate, sizeof(bitrate));

	addr.sin_family = AF_INET;
	addr.sin_port = htons(SERVER_PORT);
	addr.sin_addr.s_addr = INADDR_ANY;
	memset(addr.sin_zero, '\0', sizeof(addr.sin_zero));

	bind(sockfd, (struct sockaddr *)&addr, sizeof(addr));

	int ret = listen(sockfd, 5);

	socklen_t addr_size = sizeof(addr_s);
	newSocket = accept(sockfd, (struct sockaddr *)&addr_s, &addr_size);
	
	/* A partir deste ponto, a aplicacao pode enviar dados
	 * atraves de funcoes como send() ou sendmsg().
	 */
	 
	 return 0;
}
\end{lstlisting}

No código exemplo a porta do servidor é $2000$ e utiliza-se o código de opção de \textit{socket} \textit{GMTP\_SOCKOPT\_MEDIA\_RATE} para informar ao GMTP que o \textit{bitrate} da mídia transmitida é \ut{2}{Mbps}.

%====================================================
\chapter{Opções de \textit{Socket} do GMTP}
\label{ap:sockopt}

Neste apêndice, apresentam-se todas as opções de \textit{socket} disponíveis para utilização no GMTP, através da chamada das funções \textit{getsockopt} e \textit{setsockopt}.
Salienta-se que os protótipos das referidas funções são definidos conforme segue-se:

\begin{lstlisting}[language=C++, caption={Protótipos das funções \textit{getsockopt} e \textit{getsockopt}}, label={list:set-get-sockopt}]
int getsockopt(int s, int level, int optname, void *optval, socklen_t *optlen);
int setsockopt(int s, int level, int optname, const void *optval, socklen_t optlen);
\end{lstlisting}

Nas referidas funções, o parâmetro $s$ refere-se ao descritor do \textit{socket}, $level$ aponta se as configurações serão realizadas ao protocolo do próprio \textit{socket} ou ao seu protocolo subjacente, $optname$ refere-se ao nome da opção de \textit{socket} a ser configurada, $optval$ é um ponteiro para o valor da opção que deseja-se configurar/ler, e $optlen$ refere-se ao tamanho em \textit{bytes} de $optval$.

Para configurar ou consultar as opções de \textit{socket} do GMTP, deve-se configurar o parâmetro $level$ como o valor $281$ ($SOL\_GMTP$), conforme ilustra-se no Apêndice \ref{ap:servidor-gmtp}.
As opções de \textit{socket} disponíveis no GMTP, ordenadas pelo valor numérico de seus respectivos \textit{optnames} são:

\begin{enumerate}
	\item \textit{GMTP\_SOCKOPT\_FLOWNAME}: Configura ou consulta o nome do fluxo de mídia transmitido/recebido. Seu $optval$ é uma \textit{string} de 16 caracteres, correspondente a um código de \textit{hash} MD5 no formato UUID de 128 bits, conforme descrito na Seção \ref{sec:fluxo-basico}. Seu $optlen$ será sempre igual a 16.
	
	\item \textit{GMTP\_SOCKOPT\_MEDIA\_RATE}: Configura ou consulta o \textit{Bitrate} da mídia transmitida/recebida. Seu $optval$ é um valor inteiro não negativo de \ut{32}{bits}.
	
	\item \textit{GMTP\_SOCKOPT\_MAX\_TX\_RATE}: Configura ou consulta a taxa de transmissão máxima configurada manualmente no servidor. Seu $optval$ é um valor inteiro não negativo de \ut{32}{bits}.
	
	\item \textit{GMTP\_SOCKOPT\_UCC\_TX\_RATE}: Consulta a taxa de transmissão máxima recebida via GMTP-UCC. Essa opção está disponível apenas para leitura. Seu $optval$ é um valor inteiro não negativo de \ut{32}{bits}.
	
	\item \textit{GMTP\_SOCKOPT\_GET\_CUR\_MSS}: Retorna o MSS ou \textit{maximum segment size}. MSS é o tamanho máximo de dados que pode ser enviado, em bytes, em um único pacote. No MSS não são computados os bits pertencentes ao cabeçalho de dados. 
	Essa opção está disponível apenas para leitura. Seu $optval$ é um valor inteiro não negativo de \ut{32}{bits}.
	
	\item \textit{GMTP\_SOCKOPT\_SERVER\_RTT}: RTT (round-trip time) medido no servidor, para determinado fluxo de dados. Essa opção está disponível apenas para leitura. Seu $optval$ é um valor inteiro não negativo de \ut{32}{bits}.
	
	\item \textit{GMTP\_SOCKOPT\_SERVER\_TIMEWAIT}: \textit{Timewait} no servidor. Atualmente não é utilizado. 
	
	\item \textit{GMTP\_SOCKOPT\_PULL}: Habilita o envio de pacotes do tipo \pac{Data-Pull-Request}. Atualmente não é utilizado.
	
	\item \textit{GMTP\_SOCKOPT\_ROLE\_RELAY}: Muda ou consulta o tipo de \textit{socket}, que pode ser \textit{comum} ou \textit{repassador}. Só pode ser configurado caso o \textit{socket} esteja inativo, ou seja, não esteja no estado \textit{listening}, enviando ou recebendo dados.
	Seu $optval$ é um valor inteiro não negativo de \ut{32}{bits}, sendo igual a 1, caso o \textit{socket} seja configurado como \textit{Relay} (repassador), e 0, caso seja um \textit{socket} comum (cliente, servidor ou relator).
	
	\item \textit{GMTP\_SOCKOPT\_RELAY\_ENABLED}: Habilita ou desabilita o GMTP-Inter em determinado \textit{host}. Só é válido caso o \textit{socket} seja do tipo \textit{Relay} (repassador), ou seja, se a opção \textit{GMTP\_SOCKOPT\_ROLE\_RELAY} foi configurada com o valor 1.
	Seu $optval$ é um valor inteiro não negativo de \ut{32}{bits}, sendo igual a 1 para habilitar o GMTP-Inter e igual a 0 para desabilitar o GMTP-Inter.
	
	\item \textit{GMTP\_SOCKOPT\_UCC\_TYPE}: Configura ou consulta o modo de controle de congestionamento do GMTP-UCC. Atualmente não é utilizado. 
\end{enumerate}

\chapter{Detalhes dos experimentos}
\label{ap:experimentos}

Neste apêndice, apresentam-se detalhes a respeitos dos experimentos realizados no contexto deste trabalho.

\section{Quantidade de Ensaios}
\label{ap:ensaios}

%TODO Mudar para a equacao que foi utilizada de fato.
No experimento realizado para analisar o desempenho do GMTP, apresentado no Capítulo \ref{cap:desempenho}, fez-se necessário determinar a quantidade de repetição de um tratamento para que fosse possível analisar, com 95\% de certeza, os valores obtidos para as variáveis dependentes nos experimentos estabelecidos.
O valor de $n$ foi calculado através da Inequação~\ref{eq:n-puro}, onde $\rho^2$ e $E_0$ correspondem à variância da amostra e ao erro amostral máximo tolerado, respectivamente.

\begin{equation}
n \geq \frac{z_\gamma^2 \times \rho^2}{E_0^2}
\label{eq:n-puro}
\end{equation} 

Ao considerar o grau de confiança como $0,95$, tem-se que $E_0$ é igual a 0,05 (obtido através da subtração $1 - 0,95$) e $z_\gamma$ é igual a $1,96$, obtendo-se a Inequação~\ref{eq:n-confianca}. 

\begin{equation}
n \geq \frac{1.96^2 \times \rho^2}{0.05^2}
\label{eq:n-confianca}
\end{equation} 



\end{apendicesenv}
% ---


% ----------------------------------------------------------
% Anexos
% ----------------------------------------------------------

% ---
% Inicia os anexos
% ---
\begin{anexosenv}

% Imprime uma página indicando o início dos anexos
\partanexos

% ======================================================================
\chapter{Estruturas que representam \textit{sockets} no Linux}
\label{anex:struct_sk}

Neste anexo, apresentam-se detalhes das estruturas que representam \textit{sockets} no núcleo do Linux.

\section{Interface para o espaço de usuário}
\label{anex:socket}

No trecho de código fonte a seguir, apresenta-se parcialmente a estrutura \lstinline|struct sock|, que é a representação de \textit{sockets} no espaço de usuário, na camada de aplicação.

\begin{lstlisting}[language=C++, caption={Estrutura que representa um \textit{socket} no espaço de usuário do Linux}, label={list:socket}]
struct socket {
	socket_state			state;

	kmemcheck_bitfield_begin(type);
	short					type;
	kmemcheck_bitfield_end(type);
	...

	struct file				*file;
	struct sock				*sk;
	const struct proto_ops	*ops;
};
 
/* (include/linux/net.h) */
\end{lstlisting}

A seguir, a descrição dos membros da estrutura:
\begin{itemize}
\item \textit{state}: Um \textit{socket} pode encontrar-se em vários estados, como  \lstinline|SS_UNCONNECTED| ou \lstinline|SS_CONNECTED|, conectado e desconectado, respectivamente.
\item \textit{type}: O tipo do \textit{socket}, por exemplo, \lstinline|SOCK_STREAM| para TCP, \lstinline|SOCK_DGRAM| para UDP.
\item \textit{file}: O arquivo associado ao \textit{socket}.
\item \textit{sk}: O objeto \lstinline|sock| associado ao \textit{socket}. Representa a interface para a camada de rede.
\item \textit{ops}: Um apontador para a estrutura \lstinline|struct proto_ops|. Consiste basicamente de \textit{callbacks} para o \textit{socket}, conforme discutido na Seção \ref{sec:sockets}.
\end{itemize}

\section{Interface para a camada de rede}
\label{anex:sock}

No trecho de código fonte a seguir, apresenta-se parcialmente a estrutura \lstinline|struct sock|, que é a representação de \textit{sockets} na camada de rede.

\begin{lstlisting}[language=C++, caption={Estrutura que representa um \textit{socket} na camada de transporte do Linux}, label={list:sock}]
struct sock {
	...
	
	struct sk_buff_head	sk_receive_queue;
	...
	
	int					sk_rcvbuf;
	...

	int					sk_sndbuf;
	struct sk_buff_head	sk_write_queue;
	...
	unsigned int		sk_shutdown  : 2,
						sk_no_check_tx : 1,
						sk_no_check_rx : 1,
						...
						sk_protocol  : 8,
						sk_type      : 16;
	...
	
	void				(*sk_data_ready)(struct sock *sk);
	void				(*sk_write_space)(struct sock *sk);
	...
	
};
 
/* (include/net/sock.h) */
\end{lstlisting}

A seguir, a descrição dos membros da estrutura:
\begin{itemize}
\item \textit{sk\_receive\_queue}: Uma fila para pacotes recebidos.
\item \textit{sk\_rcvbuf}: O tamanho do \textit{buffer} de recepção em \textit{bytes}.
\item \textit{sk\_sndbuf}: O tamanho do \textit{buffer} de envio em \textit{bytes}.
\item \textit{sk\_write\_queue}: Uma fila de pacote para enviar.
\item \textit{sk\_no\_check\_tx} e \textit{sk\_no\_check\_rx}: Um marcador para desabilitar o \textit{checksum} na transmissão e na recepção, respectivamente.
\item \textit{sk\_protocol}: Um número de protocolo, conforme discutido na Seção \ref{sec:sockets}. 
\item \textit{sk\_type}: O tipo do \textit{socket}, por exemplo, \lstinline|SOCK_STREAM| para TCP, \lstinline|SOCK_DGRAM| para UDP, conforme discutido na Seção \ref{sec:sockets}.
\item \textit{sk\_data\_ready}: Um \textit{callback} para notificar o \textit{socket} sobre a chegada de dados.
\item \textit{sk\_write\_space}: Um \textit{callback} para indicar que existe memória livre disponível para transmissão de dados.
\end{itemize}


\end{anexosenv}

%---------------------------------------------------------------------
% INDICE REMISSIVO
%---------------------------------------------------------------------
%\phantompart
%\printindex
%--------------------------------------------------------------------